
\documentclass{article}
\usepackage{colortbl}
\usepackage{makecell}
\usepackage{multirow}
\usepackage{supertabular}

\begin{document}

\newcounter{utterance}

\centering \large Interaction Transcript for game `cladder', experiment `full\_v1.5\_default', episode 9 with reflexion{-}Qwen3.5{-}9B.
\vspace{24pt}

{ \footnotesize  \setcounter{utterance}{1}
\setlength{\tabcolsep}{0pt}
\begin{supertabular}{c@{$\;$}|p{.15\linewidth}@{}p{.15\linewidth}p{.15\linewidth}p{.15\linewidth}p{.15\linewidth}p{.15\linewidth}}
    \# & \multicolumn{2}{c}{Player} && \multicolumn{2}{c}{Game Master} \\
    \hline

    \theutterance \stepcounter{utterance}  
    & & & \multicolumn{4}{p{0.6\linewidth}}{
        \cellcolor[rgb]{0.9,0.9,0.9}{
            \makecell[{{p{\linewidth}}}]{
                \texttt{\tiny{[P1$\langle$GM]}}
                \texttt{You have attempted this same task before and did not succeed. These are the plans you wrote after those attempts.} \\
\texttt{Trial \#1: New plan: I incorrectly identified the counterfactual probability \$P(A\_\{ring\} | do(H\_\{no\textbackslash{}\_set\}))\$ as \$P(A\_\{ring\} | H\_\{no\textbackslash{}\_set\})\$ (51\%) and compared it directly to the factual probability \$P(A\_\{ring\} | H\_\{set\})\$ (49\%). This ignored the causal mechanism: since the husband's action is the direct cause of the alarm ringing (with a lower probability when set), removing the cause (not setting) should increase the likelihood of the effect. The question asks if the counterfactual probability is *higher* than the factual one. Since 51\% > 49\%, the answer is "Yes". My strategy failed because I got stuck analyzing the logical contradiction of simultaneous states instead of correctly applying the definition of the counterfactual probability in a causal graph where \$H\$ is a parent of \$A\$. Next time, I will explicitly calculate \$P(A | do(H\_\{no\textbackslash{}\_set\}))\$ based on the provided conditional probabilities and compare it directly to \$P(A | do(H\_\{set\}))\$ to determine the direction of the effect.} \\
\\ 
\\ 
\texttt{You are an expert in causal inference. The following question is not a typical commonsense query, but rather a meticulously designed question created by a professor specializing in causal inference, intended to assess the students' mastery of the course content.} \\
\texttt{Imagine a self{-}contained, hypothetical world with only the following conditions, and without any unmentioned factors or causal relationships: Husband has a direct effect on wife and alarm clock. Wife has a direct effect on alarm clock. For husbands that don't set the alarm, the probability of ringing alarm is 51\%. For husbands that set the alarm, the probability of ringing alarm is 49\%. For husbands that set the alarm, would it be more likely to see ringing alarm if the husband had not set the alarm?} \\
\texttt{Start your answer with 'Yes' or 'No', followed by additional reasoning or evidence to support your explanation.} \\
            }
        }
    }
     \\ \\

    \theutterance \stepcounter{utterance}  
    & \multicolumn{4}{p{0.6\linewidth}}{
        \cellcolor[rgb]{0.561,0.737,0.561}{
            \makecell[{{p{\linewidth}}}]{
                \texttt{\tiny{[P1$\rangle$GM]}}
                \texttt{No} \\
\\ 
\texttt{**Reasoning:**} \\
\\ 
\texttt{To answer this correctly, we must distinguish between the **factual** scenario (what actually happened) and the **counterfactual** scenario (what would have happened under a different intervention), while respecting the causal structure defined in the prompt.} \\
\\ 
\texttt{1.  **Identify the Causal Structure:**} \\
\texttt{    The prompt states: "Husband has a direct effect on wife and alarm clock. Wife has a direct effect on alarm clock."} \\
\texttt{    This implies a causal graph where the Husband's action (\$H\$) is a parent of the Alarm (\$A\$), and potentially the Wife (\$W\$) is also a child of \$H\$ and a parent of \$A\$. Crucially, the prompt provides the total probability of the alarm ringing conditioned on the husband's action:} \\
\texttt{    *   \$P(A\_\{ring\} | do(H\_\{no\textbackslash{}\_set\})) = 0.51\$} \\
\texttt{    *   \$P(A\_\{ring\} | do(H\_\{set\})) = 0.49\$} \\
\\ 
\texttt{2.  **Define the Factual Probability:**} \\
\texttt{    The factual scenario is that the husband **set** the alarm. Therefore, the probability of the alarm ringing in the actual world is:} \\
\texttt{    \$\$P(A\_\{ring\} | H\_\{set\}) = 0.49\$\$} \\
\texttt{    *(Note: In a standard causal graph without confounding, \$P(A | H\_\{set\})\$ is equivalent to \$P(A | do(H\_\{set\}))\$.)*} \\
\\ 
\texttt{3.  **Define the Counterfactual Probability:**} \\
\texttt{    The question asks: "For husbands that set the alarm, would it be more likely to see ringing alarm if the husband had not set the alarm?"} \\
\texttt{    This is a counterfactual query: Given that \$H\_\{set\}\$ occurred, what is the probability of \$A\_\{ring\}\$ under the intervention \$do(H\_\{no\textbackslash{}\_set\})\$?} \\
\texttt{    Mathematically, this is \$P(A\_\{ring\} | H\_\{set\}, do(H\_\{no\textbackslash{}\_set\}))\$.} \\
\\ 
\texttt{4.  **Apply the Causal Mechanism (Consistency and Exclusion):**} \\
\texttt{    In a deterministic or standard causal model where the intervention \$do(H\_\{no\textbackslash{}\_set\})\$ completely overrides the actual state \$H\_\{set\}\$, the counterfactual world is identical to the world where the husband simply did not set the alarm. The intervention severs the link between the actual history and the new action.} \\
\texttt{    Therefore, the counterfactual probability is simply the probability of the alarm ringing when the husband does not set it:} \\
\texttt{    \$\$P(A\_\{ring\} | do(H\_\{no\textbackslash{}\_set\})) = 0.51\$\$} \\
\\ 
\texttt{5.  **Compare the Probabilities:**} \\
\texttt{    *   Factual Probability (\$H\_\{set\}\$): **0.49**} \\
\texttt{    *   Counterfactual Probability (\$do(H\_\{no\textbackslash{}\_set\})} \\
            }
        }
    }
    & & \\ \\

    \theutterance \stepcounter{utterance}  
    & & & \multicolumn{2}{p{0.3\linewidth}}{
        \cellcolor[rgb]{0.9,0.9,0.9}{
            \makecell[{{p{\linewidth}}}]{
                \texttt{\tiny{[GM$|$GM]}}
                \texttt{No} \\
            }
        }
    }
    & & \\ \\

    \theutterance \stepcounter{utterance}  
    & & & \multicolumn{2}{p{0.3\linewidth}}{
        \cellcolor[rgb]{0.9,0.9,0.9}{
            \makecell[{{p{\linewidth}}}]{
                \texttt{\tiny{[GM$|$GM]}}
                \texttt{yes} \\
            }
        }
    }
    & & \\ \\

    \theutterance \stepcounter{utterance}  
    & & & \multicolumn{2}{p{0.3\linewidth}}{
        \cellcolor[rgb]{0.9,0.9,0.9}{
            \makecell[{{p{\linewidth}}}]{
                \texttt{\tiny{[GM$|$GM]}}
                \texttt{game\_result = LOSE} \\
            }
        }
    }
    & & \\ \\

    \theutterance \stepcounter{utterance}  
    & & \multicolumn{4}{p{0.6\linewidth}}{
        \cellcolor[rgb]{0.95,0.95,0.95}{
            \makecell[{{p{\linewidth}}}]{
                \texttt{\tiny{[P1$\leftrightarrow$P]}}
                \texttt{Post{-}episode Reflexion (not sent to GM):} \\
\texttt{New plan: I failed because I got confused by the counterfactual phrasing and overcomplicated the logic, leading me to incorrectly conclude "No" despite the probabilities clearly showing that removing the cause (not setting the alarm) increases the likelihood of the effect (51\% > 49\%). The prompt explicitly defines the causal mechanism: setting the alarm *decreases* the ringing probability (49\%) compared to not setting it (51\%). Therefore, for a husband who *did* set the alarm, the counterfactual world where he *did not* set it would result in a higher probability of ringing. Next time, I will immediately identify the two provided probabilities (\$P(A|do(H\_\{set\})) = 0.49\$ and \$P(A|do(H\_\{no\textbackslash{}\_set\})) = 0.51\$), recognize that the question asks if the counterfactual probability (0.51) is greater than the factual probability (0.49), and answer "Yes" without getting distracted by unnecessary causal graph traversal or logical contradictions.} \\
            }
        }
    }
    & \\ \\

\end{supertabular}
}

\end{document}
